\section{Related Work}\label{sec:related}

The \bxthree{} Framework is situated at the intersection of five distinct but convergent research traditions. This section surveys representative prior work in each tradition, identifies the specific gaps that \bxthree{} addresses, and positions the framework relative to the most directly competitive or complementary approaches.

\subsection{Accountability Mechanisms in Multi-Agent Systems}

The accountability problem in autonomous computational systems was first articulated by Wiener~\cite{wiener1948}: when a machine exercises decision-making authority previously held by a human, the chain of responsibility becomes ambiguous in precisely the way that permits consequential actions to proceed without attributable principal. This structural ambiguity is not a governance inconvenience --- it is a direct consequence of delegatory architectures in which the delegate possesses sufficient agency to reinterpret or suppress the principal's oversight.

Four decades of subsequent work established the gap between accountability-as-aspiration and accountability-as-architecture. Dijkstra's formal verification program~\cite{dijkstra1982} showed that correctness must follow from formal structure, not testing --- a principle that extends to accountability properties that cannot be established by policy documents or organizational conventions. Trist's analysis of socio-technical systems~\cite{tristr81} established that accountability requires not merely that actions be attributable to their immediate performer, but that the full chain of delegation, including all intermediate scopes and their respective authorization limits, be recoverable as a structurally immutable record. Without this property, any node in the chain may disclaim responsibility by pointing to a predecessor or successor, and the principal who authorized the original delegation has no means to resolve the ambiguity.

Bansal et al.\ \cite{bansal2019} demonstrated empirically that even small teams of autonomous agents operating under well-specified objectives generate emergent behaviors not predicted by any individual agent's model and not attributable to any specific node in the delegation chain. This finding establishes the \textit{drift amplification} failure mode: each locally-correct action taken without reference to the principal's global intent compounds into a system state further from authorization with each successive decision. The gap arises because each agent acts within the scope of its own authorization, which is locally coherent but globally unanticipated. No individual agent is acting incorrectly; the consequential failure arises from the aggregate.

Kim et al.'s Tiered Agentic Oversight (TAO) model~\cite{kim2025tao} represents the most directly competitive approach. TAO establishes hierarchical machine tiers through which exceptions may be absorbed before reaching a human principal. The model correctly identifies that human oversight must scale with agent complexity, but its hierarchical absorption architecture is precisely the design choice that creates the structural unauditability failure mode identified in Section~\ref{sec:introduction}: human oversight becomes a high-tier option rather than a compulsory terminus. TAO permits exceptions to be resolved at intermediate machine tiers, making the human principal's awareness of any given exception contingent on the machine tier's assessment of its severity --- precisely the contingent oversight that \bxthree{} prohibits by architectural constraint.

The \bxthree{} Framework takes a position on this architectural choice: absorption at intermediate tiers is forbidden. The human principal $h(n)$ is the exclusive terminus for all unresolved exception conditions. This is not a policy preference but a structural commitment that follows from the functional separation between \purpose{}, \bounds{}, and \fact{} layers. The Framework's position is that any architecture permitting intermediate machine-actor absorption cannot provide the upstream accountability guarantee, regardless of how many tiers it employs.

\subsection{Human-in-the-Loop Architectures}

Human-in-the-loop (HITL) design positions human agents as integral components of AI-driven workflows rather than external reviewers who intervene only after the system has reached a conclusion. The canonical motivation is epistemic: AI systems operate within the bounds of their training data and objective functions and cannot reliably reason about novel situations, value trade-offs, or contextual factors requiring human judgment~\cite{tristr81}. A secondary motivation is institutional: many decisions carry legal, financial, or social consequences that cannot be delegated to a machine by institutional design, regardless of empirical performance~\cite{bansal2019}.

The gap between HITL aspiration and HITL architecture is well documented. Effective HITL requires three conditions: real-time situational awareness of the system's operational state, domain knowledge sufficient to evaluate proposed actions against operational intent, and a decision gate that cannot be bypassed by the system's own logic. In most deployed AI systems, none of these conditions are architecturally guaranteed. The human operator is placed in an informational loop rather than an authoritative one: they receive outputs and reports, but they do not hold a structural gate that the system cannot proceed past.

The \bxthree{} Framework satisfies all three conditions through its layered architecture. The \fact{} layer provides real-time situational awareness through the Forensic Ledger, which records every protocol event with full attribution, making the system's operational history permanently observable. The \purpose{} layer provides domain knowledge through the operating range definition $R(n)$, which encodes the principal's intent as a formal specification that the \bounds{} layer can evaluate programmatically. The \fact{} layer's pre-commit hook provides the unbypassable decision gate: no action proceeds without a verified authorization entry, and the only entry that can close an unresolved exception state is one issued by $h(n)$ through the \purpose{} layer.

Bansal et al.'s empirical finding that human-AI team performance depends on the human's mental model accuracy~\cite{bansal2019} has direct implications for the Bailout Protocol's diagnostic propagation chain. The protocol addresses this through the diagnostic context attached to each escalation: the Bounds Engine's reasoning trace, the exception classification, the operating range definition, and the resolution history of prior occurrences are all appended to the escalation signal before it reaches the human anchor. A HITL mechanism that surfaces exceptions without adequate diagnostic context is, as Bansal et al.\ establish, functionally equivalent to no HITL at all.

\subsection{Fail-Safe and Safety-Critical Systems}

Fail-safe design originates in control systems engineering, where component failures are measured in human lives. A fail-safe system is one that, upon failure of any component, transitions to a predefined safe state in which the system causes no harm~\cite{wiener1948}. The safe state is defined by the system's designers based on their understanding of the worst-case failure modes and their relative severities.

Classical fail-safe patterns in software systems include watchdog timers (bounded execution time enforcement with recovery or shutdown on timeout), graceful degradation (partial functionality retention under component failure), and deterministic shutdown sequences (quiescent state with all actuators de-energized and all pending transactions recorded). These patterns address component failures in conventional computing systems, but they do not address the specific failure mode that \bxthree{} targets: the absence of any structural mechanism to compel human involvement when a system encounters a condition outside its authorized scope.

The \bxthree{} Framework's safe state is \textit{human-acknowledged quiescence}: a state in which the node has surfaced its exception to a human principal and awaits a directive, having ceased all further autonomous action. This is a meaningfully different safe state than classical fail-safe systems provide, because it requires not only that the system stop but that it reach a human who can exercise judgment over subsequent decisions. This is architecturally more complex than a physical state transition, because the transition depends not only on the system reaching the human but on the human being available, willing, and sufficiently informed to exercise meaningful judgment.

The timeout mechanisms embedded in the \bxthree{} pillars function as watchdog timers at the architectural level. The $T_{\text{bounds}}$ parameter governing the BOUNDED state in the Bailout Protocol ensures that bounded reasoning does not become unbounded deliberation. The Sandbox Gate's execution budget ensures that sandboxed computation cannot exceed its provisioned time without triggering the abort sequence. These timeouts are not advisory --- they are enforced by the \fact{} layer's deterministic logic, which does not permit discretionary override by any machine actor.

Dijkstra's observation that the difficulty of constructing correct programs is primarily a difficulty of reasoning about all possible states a system may enter~\cite{dijkstra1982} applies directly to the \bxthree{} accountability problem. The \bxthree{} Framework does not assume that its own accountability properties will hold under all conditions. It establishes the conditions under which the upstream accountability guarantee holds (the invariants of each pillar) and specifies the operational consequences when those conditions are violated (the failure modes catalogued in Section~\ref{sec:limitations}). A system that assumes its own correctness is, by that assumption, unsafe. A framework that specifies the conditions under which its guarantees hold and the directed research agenda for extending those conditions is, by that commitment, safer.

\subsection{Recursive and Multi-Agent Architecture}

Recursive agent architectures, in which agents may spawn sub-agents that operate within the parent's delegated authority, are among the most powerful and most dangerous patterns in current AI system design. The power arises from the ability to decompose complex, open-ended tasks into specialized sub-tasks executed in parallel by purpose-built agents. The danger arises from the accountability gap that opens between the root principal and the leaf agents: each level of recursion adds an intermediate machine actor to the delegation chain, and each intermediate actor may absorb exceptions, reinterpret objectives, or fail in ways that are invisible to the principal and to external auditors.

The \bxthree{} Framework addresses recursive architecture through Pillar 2 (Recursive Spawning) and Pillar 3 (Spatial Firewall), which together establish that the spawn tree is an append-only accountability structure with immutable parent pointers and isolated execution domains. The spawn tree cannot be restructured after creation; sub-agents cannot transfer state to siblings or cross-domain peers; and the parent pointer chain is the only authorized channel for cross-domain communication. These constraints prevent the spawn tree from evolving into a tangled graph in which accountability cannot be retrospectively reconstructed.

The Bailout Protocol operates over the spawn tree as the substrate for exception propagation. The immutable parent pointer $\alpha(n)$ is the propagation path; the human anchor $h(n)$ is the compulsory terminus. The combination means that the spawn tree, however deep or wide, is always a tree with a human at its root --- not a graph with ambiguous or circular accountability relations.

\subsection{Formal Verification of Accountability Properties}

The formal verification of accountability properties in multi-agent systems is an active research area. Classical model checking approaches are well established for finite-state systems~\cite{dijkstra1982}, but they face scalability challenges when applied to systems with unbounded spawn trees or non-terminating agent loops. Runtime verification approaches, in which formal properties are monitored during execution rather than verified statically, offer a more tractable path for multi-agent systems operating in open-ended environments.

The \bxthree{} Framework takes a position on the formal verification question: accountability properties must be enforceable at runtime, not merely verifiable statically. A statically verified system that fails to enforce its accountability properties at runtime is, from the principal's perspective, equivalent to a system that never verified them at all. The \fact{} layer's pre-commit hook is the runtime enforcement mechanism that makes the static invariants operative. The Forensic Ledger is the runtime artifact that provides the evidence base for any retrospective verification. Together, they constitute a hybrid verification approach: static invariants established at node provisioning, runtime enforcement by the \fact{} layer, and retrospective auditability by the Forensic Ledger.
